\section{Selection by a vanishing exponential parameter}
\label{sec:vanishing-parameter-selection}

The equation with exponent zero can have several normalized solutions. Nevertheless, the solutions of the equations with positive exponent select one of them in the limit: the scalar functional provides a family of integral normalizations, and the selected potential is the greatest subsolution of the critical equation satisfying all of them. To prove that it is the greatest, we use supporting probabilities of a second functional, obtained by varying the forcing term in the subsolution inequality.

In this section, $\theta$ is an arbitrary smooth big real $(1,1)$-form on $X$, not necessarily semipositive. Fix a strict analytic barrier $\rho\leq0$ as in \cref{sec:big-minimal-comparison}, and put $\Omega=X\setminus\{\rho=-\infty\}$. A potential has \emph{minimal type} if it has the ordinary singularity type of the envelope $V_\theta=P_\theta(0)$, which has minimal singularities (\cref{sec:obstacle-envelopes}). Fix a nonzero $0\leq f\in L^p(\mathrm{d}V)$, with $1<p<\infty$. For a nonzero nonnegative density $g$ in some $L^q$, $q>1$, let $c(g)>0$ be the unique constant of \cref{thm:big-minimal-normalized} for which the equation $\mathrm{MA}_\theta(v)=c(g)g\,\mathrm{d}V$ has a solution $v$ of minimal type. We call $c(g)$ the \emph{critical scalar} of $g$ and such solutions \emph{critical potentials} for $g$; they are continuous on $\Omega$ by \cref{thm:big-minimal-holder}. Here $\mathrm{MA}_\theta$ is the extended measure $\mathrm{MA}^{\mathrm{ext}}_\theta$. On $\Omega$ it is the locally bounded product for every potential of minimal type, and it gives zero mass to the analytic set $X\setminus\Omega$.

Put $c=c(f)$ and define
\begin{equation}\label{eq:selection-probabilities}
 \begin{aligned}
 \kappa(r)&=\log c-\log c(\mathrm{e}^r f),\\
 \mathcal P_f&=\left\{\nu\in\mathcal P(X):
       \int_X r\,\mathrm{d}\nu\leq\kappa(r)
       \text{ for every }r\in C(X)\right\}.
 \end{aligned}
\end{equation}
When $\theta$ is semipositive and $f=\mathrm{e}^F$ with $F\in C(X)$, we have $\kappa(r)=B_0(F)-B_0(F+r)$, and $\mathcal P_f$ is the set $\mathcal A_0(F)$ of supporting probabilities. The definition also applies when $f$ vanishes on a set of positive measure and when $\theta$ is not semipositive.

\begin{theorem}\label{thm:vanishing-exponential-selection}\index{Monge--Ampere equation@Monge--Amp\`ere equation!vanishing-exponent selection}
For $\lambda>0$, let $u_\lambda$ be the solution of minimal type of
\[
 \mathrm{MA}_\theta(u_\lambda)
       =\mathrm{e}^{\lambda u_\lambda}f\,\mathrm{d}V
\]
given by \cref{thm:big-minimal-exponential} and unique by \cref{cor:big-barrier-orders}. The set $\mathcal P_f$ is nonempty and weakly compact, and every potential of minimal type is integrable against all its measures. As $\lambda\downarrow0$,
\[
 w_\lambda=u_\lambda-\lambda^{-1}\log c
\]
converges locally uniformly on $\Omega$ and in $L^1(X)$ to a potential $w$ of minimal type, continuous on $\Omega$. It is the greatest among the potentials $s$ of minimal type, continuous on $\Omega$, satisfying
\begin{equation}\label{eq:selection-class}
 \mathrm{MA}_\theta(s)\geq cf\,\mathrm{d}V,\qquad
 \int_Xs\,\mathrm{d}\nu\leq0\quad(\nu\in\mathcal P_f).
\end{equation}
Moreover,
\[
 \mathrm{MA}_\theta(w)=cf\,\mathrm{d}V,\qquad
 \max_{\nu\in\mathcal P_f}\int_Xw\,\mathrm{d}\nu=0.
\]
Consequently, $u_\lambda-\sup_Xu_\lambda$ converges to $w-\sup_Xw$ in the same two topologies. For a Hermitian $\theta$, one may take $\Omega=X$, and both convergences of potentials are uniform.
\end{theorem}

\subsection{A norm defined by the critical scalar}

Mixed products turn the comparison of critical scalars into a triangle inequality for $n$th roots of densities, and the subsolution estimate associated with the strict barrier gives continuity in the corresponding Lebesgue norm.

Fix $1<q<p$ and set $\ell=nq$. On real $L^\ell(\mathrm{d}V)$ define
\[
 N(a)=
 \begin{cases}
 c(|a|^n)^{-1/n},&a\ne0,\\
 0,&a=0.
 \end{cases}
\]

\Needspace{12\baselineskip}
\begin{proposition}\label{prop:critical-scalar-norm}
The function $N$ is a lattice norm on $L^\ell(\mathrm{d}V)$, that is, a norm satisfying $N(a)\leq N(b)$ whenever $|a|\leq|b|$. There is $C_q<\infty$, depending only on $X,\theta,\mathrm{d}V$ and $q$, such that
\begin{equation}\label{eq:critical-scalar-norm-bound}
 N(a)\leq C_q\|a\|_\ell,\qquad
 |N(a)-N(b)|\leq C_q\|a-b\|_\ell.
\end{equation}
Consequently, for nonzero $g,h\geq0$ in $L^q$,
\[
 |c(g)^{-1/n}-c(h)^{-1/n}|
       \leq C_q\|g-h\|_q^{1/n}.
\]
\end{proposition}
\begin{proof}
If a potential $v$ of minimal type satisfies $\mathrm{MA}_\theta(v)\geq d g\,\mathrm{d}V$, then $d\leq c(g)$, since otherwise comparison with a critical potential for $g$ contradicts part (1) of \cref{cor:big-barrier-orders}. Hence $g\leq h$ implies $c(g)\geq c(h)$. Uniqueness of the critical scalar also gives $c(tg)=c(g)/t$ for $t>0$. These facts prove that $N$ is monotone and absolutely homogeneous, and positivity follows from $0<c(g)<\infty$.

For the triangle inequality, let $a,b\geq0$ be nonzero, put $A=N(a)$ and $B=N(b)$, and choose critical potentials $v_a,v_b$ with measures $(a/A)^n\mathrm{d}V$ and $(b/B)^n\mathrm{d}V$. On every relatively compact subset of $\Omega$, \eqref{eq:interior-minkowski-background} gives
\[
 \mathrm{MA}_\theta\left(\frac{A v_a+B v_b}{A+B}\right)
       \geq\left(\frac{a+b}{A+B}\right)^n\mathrm{d}V,
\]
where both potentials are locally bounded. Exhausting $\Omega$ gives the inequality on $X$, since $X\setminus\Omega$ has zero smooth volume. The averaged potential has minimal type, so the first paragraph gives $N(a+b)\leq A+B$. The cases where $a$ or $b$ vanishes are trivial, and for real functions we use $|a+b|\leq|a|+|b|$ and monotonicity.

Next, apply \cref{lem:big-reference-subsolution} with $\eta=V_\theta$ to nonnegative densities of $L^q$ norm at most one. It gives a uniform $m>0$ and subsolutions of minimal type, hence $c(g)\geq m$ for each such nonzero density. For $a\ne0$, take $g=|a|^n/\|a\|_\ell^n$, which has $L^q$ norm one. Homogeneity gives $N(a)\leq m^{-1/n}\|a\|_\ell$, which is the first bound in \eqref{eq:critical-scalar-norm-bound}, and the reverse triangle inequality gives the second. Applying the second bound to $g^{1/n},h^{1/n}$ and using $|g^{1/n}-h^{1/n}|^n\leq|g-h|$ proves the last assertion.
\end{proof}

\subsection{Norming probabilities and integral bounds}

Set $a=f^{1/n}\in L^\ell(\mathrm{d}V)$ and $A=N(a)=c^{-1/n}$. The probabilities in \eqref{eq:selection-probabilities} come from positive linear functionals that norm $N$ at $a$, that is, functionals dominated by $N$ and equal to $N(a)$ at $a$. This representation will control integrals of unbounded potentials as well as bounded logarithmic perturbations.

With $\ell'=\ell/(\ell-1)$, define
\begin{equation}\label{eq:selection-norming-face}
 \mathcal H_a=\left\{\eta\in L^{\ell'}(\mathrm{d}V):
 \eta\geq0,\ \int_Xa\eta\,\mathrm{d}V=A,\
 \int_Xb\eta\,\mathrm{d}V\leq N(b)\text{ for every }b\in L^\ell\right\}.
\end{equation}

\Needspace{18\baselineskip}
\begin{proposition}\label{prop:selection-norming-probabilities}
The set $\mathcal H_a$ is nonempty and weakly compact in $L^{\ell'}$. The function $\kappa$ of \eqref{eq:selection-probabilities} is given by
\[
 \kappa(r)=n\log\frac{N(a\mathrm{e}^{r/n})}{A},
\]
and it is convex, increasing and $1$-Lipschitz for the uniform norm on $C(X)$, with $\kappa(0)=0$ and $\kappa(r+b)=\kappa(r)+b$ for constants $b$. Its subdifferential at zero is
\begin{equation}\label{eq:selection-face-representation}
 \mathcal P_f=\{\nu_\eta:\eta\in\mathcal H_a\},\qquad
 \mathrm{d}\nu_\eta=A^{-1}a\eta\,\mathrm{d}V.
\end{equation}
If $z$ is measurable and $az\in L^\ell$, then
\begin{equation}\label{eq:selection-weighted-moment}
 \sup_{\nu\in\mathcal P_f}\int_X|z|\,\mathrm{d}\nu
       \leq C_q A^{-1}\|az\|_\ell.
\end{equation}
If also $a\mathrm{e}^{tz/n}\in L^\ell$ and $N(a\mathrm{e}^{tz/n})=A$ for some $t>0$, then $\int_Xz\,\mathrm{d}\nu\leq0$ for every $\nu\in\mathcal P_f$.
\end{proposition}
\begin{proof}
The formula for $\kappa$ follows from $N(a\mathrm{e}^{r/n})=c(\mathrm{e}^rf)^{-1/n}$. The Hahn--Banach theorem gives a linear functional that norms $N$ at $a$. By \eqref{eq:critical-scalar-norm-bound} it is continuous on $L^\ell$, and hence is represented by some $\eta_0\in L^{\ell'}$. For $b\geq0$, testing on $\operatorname{sgn}(\eta_0)b$ gives
\[
 \int_Xb|\eta_0|\,\mathrm{d}V\leq N(b).
\]
Since $a\geq0$ and $\int_Xa\eta_0\,\mathrm{d}V=A=N(a)$, the same equality holds with $|\eta_0|$ in place of $\eta_0$. Monotonicity of $N$ extends the domination to real $b$, so $|\eta_0|\in\mathcal H_a$. The defining conditions of $\mathcal H_a$ are weakly closed and give $\|\eta\|_{\ell'}\leq C_q$, so reflexivity proves weak compactness.

For nonzero nonnegative $b_0,b_1$, the weighted arithmetic--geometric mean inequality applied to $b_j/N(b_j)$, together with monotonicity and the triangle inequality, gives
\[
 N(b_0^{1-t}b_1^t)\leq N(b_0)^{1-t}N(b_1)^t,\qquad 0\leq t\leq1.
\]
Taking $b_j=a\mathrm{e}^{r_j/n}$ proves convexity of $\kappa$. Monotonicity and the translation rule follow from the corresponding properties of $N$, and together they imply the Lipschitz bound.

\Needspace{5\baselineskip}
Every $\nu_\eta$ is a probability and satisfies
\[
 \frac{N(a\mathrm{e}^{r/n})}{A}
       \geq\int_X\mathrm{e}^{r/n}\,\mathrm{d}\nu_\eta.
\]
Jensen's inequality gives $\kappa(r)\geq\int_Xr\,\mathrm{d}\nu_\eta$, which proves one inclusion in \eqref{eq:selection-face-representation}. The map $\eta\mapsto\nu_\eta$ is weakly continuous, since $ar\in L^\ell$ for $r\in C(X)$.

Conversely, let $\nu\in\mathcal P_f$ and put $L(r)=\int_Xr\,\mathrm{d}\nu$. Differentiating the curve $t\mapsto a\mathrm{e}^{tr/n}$ in $L^\ell$ at $t=0$ gives
\[
 N(a\mathrm{e}^{tr/n})
       =N(a)+\frac{t}{n}N'_+(a;ar)+o(t),\qquad t\downarrow0.
\]
Here the one-sided derivative $N'_+(a;\cdot)$ of the continuous convex function $N$ exists, and the error estimate follows from its Lipschitz bound. The inequality $L\leq\kappa$, applied to $tr$ and $-tr$, therefore gives
\begin{equation}\label{eq:selection-face-extension}
 |A L(r)|\leq C_q\|ar\|_\ell,\qquad
 A L(r)\leq N'_+(a;ar)\leq N(a+ar)-A.
\end{equation}
Thus $T(ar)=A L(r)$ defines a continuous linear functional on the space of products $ar$, $r\in C(X)$. These products are dense in
\[
 S_a=\{b\in L^\ell:b=0\text{ almost everywhere on }\{a=0\}\}:
\]
indeed, it suffices to approximate $b/a$ in $L^\ell(a^\ell\mathrm{d}V)$ by continuous functions. If $ar\geq0$, then $ar=ar_+$ almost everywhere, and the first inequality in \eqref{eq:selection-face-extension} gives $L(r-r_+)=0$, so $T(ar)\geq0$. Approximating nonnegative elements of $S_a$ by products $ar$ with $r\geq0$ continuous shows that the continuous extension of $T$ to $S_a$ is positive.

Represent this extension by some $\eta\geq0$ in $L^{\ell'}$ vanishing on $\{a=0\}$. Since $T(a)=A$, extending the second inequality in \eqref{eq:selection-face-extension} to $S_a$ and applying it to $b-a$ gives $T(b)\leq N(b)$ for $b\in S_a$. For general $b\in L^\ell$,
\[
 \int_Xb\eta\,\mathrm{d}V
 \leq N(b\mathbf1_{\{a>0\}})\leq N(b).
\]
Thus $\eta\in\mathcal H_a$, and $\nu_\eta$ agrees with $\nu$ on all continuous functions. This proves the representation \eqref{eq:selection-face-representation}, and weak compactness of $\mathcal P_f$ follows from that of $\mathcal H_a$.

The bound \eqref{eq:selection-weighted-moment} follows from H\"older's inequality and $\|\eta\|_{\ell'}\leq C_q$. Under the last hypothesis, the norming inequality gives $\int_X\mathrm{e}^{tz/n}\,\mathrm{d}\nu\leq1$. Since $z$ is $\nu$-integrable by \eqref{eq:selection-weighted-moment}, Jensen's inequality implies $\int_Xz\,\mathrm{d}\nu\leq0$.
\end{proof}

These estimates apply uniformly to all $v\in\PSH(X,\theta)$ with $\sup_Xv=0$. Indeed, \cref{lem:uniform-exponential}, applied in a Hermitian background larger than $\theta$, gives uniform bounds for all finite moments of $|v|$, and H\"older's inequality gives
\begin{equation}\label{eq:selection-psh-moments}
 \int_X |v|^{nq} f^q\,\mathrm{d}V
 \leq \|f\|_p^q
       \left(\int_X|v|^{nqp/(p-q)}\,\mathrm{d}V\right)^{(p-q)/p}.
\end{equation}
Consequently, \eqref{eq:selection-weighted-moment} gives a uniform bound for $\int_X|v|\,\mathrm{d}\nu$ over all such $v$ and all $\nu\in\mathcal P_f$; constant translates of such $v$ are integrable as well. The representation also shows that these probabilities charge neither $\{f=0\}$ nor $X\setminus\Omega$.

\subsection{Bounds for the centered potentials}

We call the potentials $w_\lambda$ of \cref{thm:vanishing-exponential-selection} the \emph{centered potentials}. Choose a critical potential $\psi$ for $f$ with $\sup_X\psi=0$, so that $\mathrm{MA}_\theta(\psi)=cf\,\mathrm{d}V$. The integral normalizations prevent the centered potentials from drifting upward by constants.

\Needspace{10\baselineskip}
\begin{lemma}\label{lem:selection-centered-bounds}
There is $K<\infty$, independent of $0<\lambda\leq1$, such that
\begin{equation}\label{eq:selection-centered-bounds}
 \psi\leq w_\lambda,\qquad
 0\leq k_\lambda:=\sup_Xw_\lambda\leq K,\qquad
 \psi-K\leq v_\lambda:=w_\lambda-k_\lambda\leq V_\theta.
\end{equation}
Moreover $\int_Xw_\lambda\,\mathrm{d}\nu\leq0$ for every $\nu\in\mathcal P_f$. The family $w_\lambda$ is relatively compact for locally uniform convergence on $\Omega$, and every subsequence converging locally uniformly on $\Omega$ also converges in $L^1(X)$, to a potential of minimal type.
\end{lemma}
\begin{proof}
The centered potential satisfies $\mathrm{MA}_\theta(w_\lambda)=c\mathrm{e}^{\lambda w_\lambda}f\,\mathrm{d}V$, while $\psi\leq0$ gives $\mathrm{MA}_\theta(\psi)\geq c\mathrm{e}^{\lambda\psi}f\,\mathrm{d}V$. Part (2) of \cref{cor:big-barrier-orders} therefore gives $\psi\leq w_\lambda$, hence $k_\lambda\geq0$. For each fixed $\lambda$, the density $\mathrm{e}^{\lambda w_\lambda}f$ lies in $L^p$, and its critical scalar is $c$. Therefore
\[
 N(a\mathrm{e}^{\lambda w_\lambda/n})=A.
\]
The moment estimate above gives $aw_\lambda\in L^\ell$, so \cref{prop:selection-norming-probabilities}, with $z=w_\lambda$ and $t=\lambda$, gives the asserted integral inequalities. Fix one $\nu\in\mathcal P_f$. Then
\[
 k_\lambda\leq-\int_Xv_\lambda\,\mathrm{d}\nu
       \leq C_q A^{-1}\|a v_\lambda\|_\ell\leq K
\]
by \eqref{eq:selection-psh-moments}, with $K$ independent of $\lambda$. The normalization $\sup_Xv_\lambda=0$ gives $v_\lambda\leq V_\theta$, and subtracting $k_\lambda$ from the lower bound $\psi\leq w_\lambda$ proves \eqref{eq:selection-centered-bounds}.

The potentials $v_\lambda=u_\lambda-\sup_Xu_\lambda$ form the normalized exponential family of \cref{prop:big-exponential-clusters}, which gives subsequences converging locally uniformly on $\Omega$ and in $L^1(X)$ to potentials of minimal type. Extracting a convergent subsequence of $k_\lambda\in[0,K]$ gives the same compactness for $w_\lambda=v_\lambda+k_\lambda$. The bound $|w_\lambda|\leq|\psi|+K$ also gives $L^1$ convergence along every subsequence converging locally uniformly on $\Omega$, since $X\setminus\Omega$ has zero volume.
\end{proof}

\Needspace{4\baselineskip}
In particular,
\[
 \mathrm{e}^{\lambda\sup_Xu_\lambda}
       =c\mathrm{e}^{\lambda k_\lambda}\longrightarrow c.
\]
\Needspace{10\baselineskip}
Moreover, for $0<\lambda\leq1$ the densities of the equations satisfy $\|c\mathrm{e}^{\lambda w_\lambda}f\|_p\leq c\mathrm{e}^K\|f\|_p$. It matters that this bound includes the scalar factor: an $L^p$ bound on the data alone, after an arbitrary rescaling, does not bound the normalized potentials. For example, on $\mathbb P^1$, with $\omega_{\mathrm{FS}}=\ddc\log(1+|z|^2)$ and $0<\epsilon<1$, the function
\[
 \psi_\epsilon=\log\frac{|z|^2+\epsilon^2}{1+|z|^2}
\]
is smooth and $\omega_{\mathrm{FS}}$-plurisubharmonic, with supremum zero and infimum $\log\epsilon^2$. Write $\mathrm{MA}_{\omega_{\mathrm{FS}}}(\psi_\epsilon)=g_\epsilon\omega_{\mathrm{FS}}$, and take $L^p$ norms with respect to $\omega_{\mathrm{FS}}$ in this example. The data $g_\epsilon/\|g_\epsilon\|_p$ have norm one, and $\psi_\epsilon$ solves the equation with these data and scalar $\|g_\epsilon\|_p$; but the densities $g_\epsilon$ of the equations themselves are not bounded in $L^p$, and the potentials are not uniformly bounded.

\subsection{Scalar envelopes of subsolutions}

To compare a cluster potential with every member of the class \eqref{eq:selection-class}, fix a point $x\in\Omega$ and add a continuous forcing term $r$ to the subsolution inequality. Let $\mathcal S$ be the convex set of potentials of minimal type that are continuous on $\Omega$, and define
\begin{equation}\label{eq:selection-scalar-envelope}
 \mathcal F_{\lambda,x}(r)=
 \sup\{s(x):s\in\mathcal S,\
       \mathrm{MA}_\theta(s)\geq
              c\mathrm{e}^{\lambda s+r}f\,\mathrm{d}V\}.
\end{equation}

\begin{lemma}\label{lem:selection-forcing-probability}
The functional $\mathcal F_{\lambda,x}$ is finite, concave, decreasing and $\lambda^{-1}$-Lipschitz for the uniform norm on $C(X)$. It satisfies
\[
 \mathcal F_{\lambda,x}(0)=w_\lambda(x),\qquad
 \mathcal F_{\lambda,x}(r+b)
       =\mathcal F_{\lambda,x}(r)-b/\lambda.
\]
There is a probability measure $\nu_{\lambda,x}$ such that, for every $r\in C(X)$ and every $s$ admissible for $r$ in \eqref{eq:selection-scalar-envelope},
\begin{equation}\label{eq:selection-forcing-support}
 \lambda s(x)+\int_Xr\,\mathrm{d}\nu_{\lambda,x}
       \leq\lambda w_\lambda(x).
\end{equation}
\end{lemma}
\begin{proof}
By part (2) of \cref{cor:big-barrier-orders}, $w_\lambda$ is the greatest admissible potential for $r=0$. If $m=\min_Xr$ and $M=\max_Xr$, then $w_\lambda-M/\lambda$ is admissible for $r$, and translating any potential admissible for $r$ by $m/\lambda$ makes it admissible for $r=0$. Hence
\[
 w_\lambda(x)-M/\lambda\leq\mathcal F_{\lambda,x}(r)
       \leq w_\lambda(x)-m/\lambda.
\]
The mixed density inequality on $\Omega$ and the arithmetic--geometric mean inequality show that a convex combination of potentials admissible for $r_0$ and $r_1$ is admissible for the same convex combination of the forcing terms. Approximating the two suprema at $x$ proves concavity. Inclusion of the admissible classes proves monotonicity, and translation by constants proves the translation rule and the Lipschitz bound.

Now apply the Hahn--Banach construction in the proof of \cref{thm:weak-adjoint-probabilities} to $-\lambda\mathcal F_{\lambda,x}$ at zero. This functional is finite, continuous, convex and increasing, and it satisfies $-\lambda\mathcal F_{\lambda,x}(r+b)=-\lambda\mathcal F_{\lambda,x}(r)+b$ for constants $b$. The same tests with nonnegative functions and constants show that its subgradient is a probability measure $\nu_{\lambda,x}$. The resulting supporting inequality, combined with $s(x)\leq\mathcal F_{\lambda,x}(r)$, gives \eqref{eq:selection-forcing-support}; it is not necessary that the supremum be attained.
\end{proof}

\begin{proof}[of \cref{thm:vanishing-exponential-selection}]
By \cref{prop:selection-norming-probabilities} and the moment bound following it, $\mathcal P_f$ has the stated properties, and \cref{lem:selection-centered-bounds} gives precompactness of the centered potentials. It remains to identify every cluster potential.

\emph{Step 1: limits of the forcing probabilities.} Fix $x\in\Omega$. For $r\in C(X)$, choose a critical potential $\psi_r$ for $\mathrm{e}^r f$ with supremum zero. By the definition of $\kappa$,
\[
 \mathrm{MA}_\theta(\psi_r)
       =c\mathrm{e}^{-\kappa(r)+r}f\,\mathrm{d}V.
\]
Since $\psi_r\leq0$, the potential $\psi_r-\kappa(r)/\lambda$ is admissible in \eqref{eq:selection-scalar-envelope}. Thus
\begin{equation}\label{eq:selection-forcing-limit}
 \int_Xr\,\mathrm{d}\nu_{\lambda,x}
       \leq\kappa(r)+\lambda\bigl(w_\lambda(x)-\psi_r(x)\bigr).
\end{equation}
For fixed $x$ and $r$, the values $w_\lambda(x)$ are uniformly bounded by \cref{lem:selection-centered-bounds}, and $\psi_r(x)$ is finite. Hence every weak cluster point of $\nu_{\lambda,x}$ as $\lambda\downarrow0$ belongs to $\mathcal P_f$.

\emph{Step 2: a pointwise comparison.} Let $s$ satisfy \eqref{eq:selection-class}. The function $-\lambda s$ is lower semicontinuous and bounded below on $X$, with the value $+\infty$ allowed, so metric inf-convolution and truncation give continuous functions $r_j\uparrow-\lambda s$. Since $r_j\leq-\lambda s$, the potential $s$ is admissible for every $r_j$, and \eqref{eq:selection-forcing-support} bounds the integrals of the $r_j$ from above. Monotone convergence, after subtracting a common lower bound, shows that $s$ is integrable against $\nu_{\lambda,x}$ and gives
\begin{equation}\label{eq:selection-pointwise-comparison}
 w_\lambda(x)\geq s(x)-\int_Xs\,\mathrm{d}\nu_{\lambda,x}.
\end{equation}

\emph{Step 3: cluster potentials belong to the class \eqref{eq:selection-class}.} Take a sequence $\lambda_j\downarrow0$ with $w_{\lambda_j}\to z$ locally uniformly on $\Omega$ and in $L^1(X)$. The bounds \eqref{eq:selection-centered-bounds} show that $z$ has minimal type, and $z$ is continuous on $\Omega$. Since $c\mathrm{e}^{\lambda_j w_{\lambda_j}}f\to cf$ in $L^p$, by dominated convergence and the common bound $w_\lambda\leq K$, continuity of local bounded products gives $\mathrm{MA}_\theta(z)=cf\,\mathrm{d}V$ on $\Omega$, and then on $X$ because neither side charges $X\setminus\Omega$.

For each $\nu\in\mathcal P_f$, the bound $|w_\lambda|\leq|\psi|+K$ has a $\nu$-integrable right-hand side, and $\nu$ charges only $\Omega$. Dominated convergence in $\int_Xw_\lambda\,\mathrm{d}\nu\leq0$ therefore gives $\int_Xz\,\mathrm{d}\nu\leq0$. Hence $z$ belongs to the class \eqref{eq:selection-class}.

\emph{Step 4: $z$ is the greatest element of the class.} For any $s$ in the class and fixed $x\in\Omega$, pass to a further subsequence with $\nu_{\lambda_j,x}\rightharpoonup\nu$, so that $\nu\in\mathcal P_f$ by Step 1. The potential $s$ is upper semicontinuous and bounded above on $X$, so
\[
 \limsup_j\int_Xs\,\mathrm{d}\nu_{\lambda_j,x}
       \leq\int_Xs\,\mathrm{d}\nu\leq0.
\]
For an unbounded $s$, this follows from the portmanteau theorem applied to $\max(s,-M)$, followed by $M\to\infty$; the final integral is finite by \eqref{eq:selection-weighted-moment}. Passing to the limit in \eqref{eq:selection-pointwise-comparison} gives $z(x)\geq s(x)$. This holds for every $x\in\Omega$, and hence on $X$ by \cref{lem:qpsh-almost-everywhere-order}. Thus every cluster potential is the greatest element $w$ of the class, and precompactness gives convergence of the whole family.

\emph{Step 5: the normalizations.} The maximum of $\int_Xw\,\mathrm{d}\nu$ over $\mathcal P_f$ is attained, because under \eqref{eq:selection-face-representation} it is the maximum of the weakly continuous functional $\eta\mapsto A^{-1}\int_Xaw\eta\,\mathrm{d}V$ on $\mathcal H_a$; note that $aw\in L^\ell$. If the maximum were negative, a small positive constant translate of $w$ would still satisfy \eqref{eq:selection-class}, contrary to greatestness. Thus it is zero.

Finally, $k_\lambda\to\sup_Xw$. Indeed, from any subsequence choose a further one with $k_\lambda\to k$. Then $v_\lambda=w_\lambda-k_\lambda$ converges in $L^1$ to $w-k$, and \cref{thm:normalized-compactness} gives $\sup_X(w-k)=0$. Thus $k=\sup_Xw$ for every cluster point of $(k_\lambda)$. Subtracting the constants proves the convergence of $u_\lambda-\sup_Xu_\lambda$. If $\theta$ is Hermitian, we may take $\rho=0$ and $\Omega=X$, so that local uniform convergence is uniform convergence.
\end{proof}

Together with \cref{prop:big-exponential-clusters}, the theorem also gives convergence of the normalized potentials in the topology of uniform convergence of differences on $\Omega$ and in every local $C^\beta$ topology stated there: compactness in those topologies and uniqueness of the $L^1$ limit force every subsequence to have the same limit. The fact that the limit is selected does not imply that the critical equation has a unique normalized solution.
